% Generated draft: requires mathematical and source review.
\documentclass[11pt]{article}
\usepackage[margin=25mm]{geometry}
\usepackage{fontspec}
\setmainfont{Latin Modern Roman}
\usepackage{amsmath,amssymb,mathtools,unicode-math}
\setmathfont{Latin Modern Math}
\usepackage{enumitem,needspace,xurl,hyperref}
\setlength{\parindent}{0pt}
\setlength{\parskip}{5pt}
\setlength{\emergencystretch}{3em}
\newcommand{\N}{\mathbb{N}}
\newcommand{\Z}{\mathbb{Z}}
\newcommand{\Q}{\mathbb{Q}}
\newcommand{\R}{\mathbb{R}}
\newcommand{\C}{\mathbb{C}}
\newcommand{\F}{\mathbb{F}}
\newcommand{\A}{\mathbb{A}}
\newcommand{\PP}{\mathbb{P}}
\newcommand{\E}{\mathbb{E}}
\newcommand{\eps}{\varepsilon}
\newcommand{\dd}{\,\mathrm{d}}
\begin{document}
\section*{UnsolvedMath ID: 10000071}
\textbf{Source problem code:} catalogue entry
\subsection{Definitions and mathematical statement}
For an integer dimension \(d\ge1\), a partition of \(\R^d\) into curve traces would consist of an index set \(I\), intervals \(T_i\), and continuous maps \(\gamma_i:T_i\to\R^d\) whose images \(L_i=\gamma_i(T_i)\) satisfy
\[
 \bigcup_{i\in I}L_i=\R^d,\qquad L_i\cap L_j=\varnothing\quad(i\ne j).
\]
The source asks for which dimensions such a partition can be made with every leaf having ``the law or geometric regularity of a Brownian path,'' and requests a rigorous construction or an impossibility proof. Standard \(d\)-dimensional Brownian motion started at \(x\) is a continuous process \(B\) with \(B_0=x\), independent increments, and \(B_t-B_s\) centered Gaussian with covariance \((t-s)I_d\) for \(0\le s<t\). This definition alone does not specify the question: a Brownian law concerns random parametrized paths, while a deterministic trace may share only selected geometric properties with those paths.

A precise target requires the source to choose a time domain, a probability space if a law is intended, and the exact Brownian property required of the leaves. In a random-partition version, the displayed partition requirements would normally be required almost surely; this quantifier is conditional on that interpretation. Whether leaves must be injectively parametrized, whether they are one-sided or two-sided, how a leaf is rooted and sampled when its law is specified, and whether any translation invariance or measurable structure is required are all absent. Thus the dimensional classification cannot be uniquely formalized from the supplied alternative ``law or geometric regularity.
\subsection{Short English statement}
The question asks which Euclidean dimensions admit a partition into disjoint Brownian-type curve traces. The source record does not define Brownian-type, and equality in distribution and resemblance of geometric properties give different requirements. The supplied open-triage assessment reports no clean resolution for all dimensions; this is not an independent status determination.
\subsection{Sources}
\begin{itemize}
\item Local UnsolvedMath registry, permanent ID \(10000071\), exact source code \texttt{catalogue entry}, ``Foliations of Euclidean space by Brownian paths''; attributed to Itai Benjamini, proposed year 2012.
\item Benjamini, \emph{Euclidean vs. Graph Metric}, Section 4.1, unnumbered Brownian-foliation question, PDF page 5: \url{https://www.wisdom.weizmann.ac.il/~itai/erd100.pdf#page=5}. Preserved source: \url{https://arquivo.pt/noFrame/replay/20201231041538id_/http://www.wisdom.weizmann.ac.il/~itai/erd100.pdf}.
\item Supplied verification dated 2026-07-29: pypdf text checked against rendered PDF; cross-document deduplication and source-side status review. The source labels the item open or conjectural, records no resolution, and the bounded modern audit found no complete primary-literature resolution. The summary mentions space-filling curves and Hausdorff-dimension-2 sets without giving a theorem or further citation.
\item Dataset attribution: Ulam AI / UnsolvedMath Contributors, CC BY 4.0. Source rights NEEDS\_REVIEW: publicly accessible author-hosted or publisher copy, with release reuse terms requiring review.
\end{itemize}
\end{document}
